\documentclass[11pt]{article}
\usepackage[margin=1in]{geometry}
\usepackage{enumitem}
% =============================================================================
% Preambles/header.tex — shared LaTeX/Beamer preamble for this project
%
% Use from a Beamer lecture by prepending:
%     \input{header}
% and compiling with TEXINPUTS=../Preambles:$TEXINPUTS (the /compile-latex
% skill does this automatically).
%
% PALETTE CONTRACT. Color names defined here MUST match the names in
% Quarto/theme-template.scss so Beamer and Quarto renderings use the same
% palette. The scripts/check-palette-sync.sh script enforces this.
%
% To customize: edit the HEX values below AND the matching $variable in
% Quarto/theme-template.scss, then run scripts/check-palette-sync.sh.
% =============================================================================

% -----------------------------------------------------------------------------
% Engine detection + encoding
%
% Our /compile-latex skill uses XeLaTeX, where inputenc is unnecessary and
% can produce encoding warnings. Only load inputenc under pdfTeX.
% -----------------------------------------------------------------------------
\usepackage{iftex}
\ifPDFTeX
  \usepackage[utf8]{inputenc}
\fi

\usepackage{amsmath, amssymb, amsthm}
\usepackage{booktabs}
\usepackage{graphicx}

% -----------------------------------------------------------------------------
% PALETTE — mirrors Quarto/theme-template.scss variable names.
% Load xcolor and define colors BEFORE anything that references them
% (notably hyperref, hypersetup, and the Beamer color assignments).
% -----------------------------------------------------------------------------
\usepackage{xcolor}

\definecolor{primary-blue}{HTML}{012169}      % headings, accents
\definecolor{primary-gold}{HTML}{B9975B}      % emphasis, borders
\definecolor{highlight-yellow}{HTML}{F2A900}  % markers, alerts
\definecolor{light-bg}{HTML}{E8EDF5}          % title / section backgrounds
\definecolor{jet}{HTML}{1A1A1A}               % body text

% Semantic colors (used in diagrams and annotations)
\definecolor{positive}{HTML}{15803D}          % good / observed / identified
\definecolor{negative}{HTML}{B91C1C}          % bad / problematic / bias
\definecolor{neutral}{HTML}{525252}           % reference / context

% Highlight accents (match .hi-* classes in the SCSS)
\definecolor{hi-slate}{HTML}{314F4F}
\definecolor{hi-green}{HTML}{2E8B57}
\definecolor{hi-red}{HTML}{C41E3A}

% Semantic aliases so skill templates and snippets can write
% \textcolor{accent}{...} without hard-coding "primary-blue".
\colorlet{accent}{primary-blue}
\colorlet{accent2}{primary-gold}

% -----------------------------------------------------------------------------
% Hyperref — loaded AFTER xcolor + palette so link colors resolve.
% -----------------------------------------------------------------------------
\usepackage{hyperref}
\hypersetup{colorlinks=true, linkcolor=primary-blue, urlcolor=primary-blue,
            citecolor=primary-blue, pdfborder={0 0 0}}

% -----------------------------------------------------------------------------
% Course code — set once per deck (after \input{header}, before \begin{document})
% via \coursecode{CS401}. Shown in the footer on every slide so a deck's course
% is identifiable without opening the title page. Empty by default.
% -----------------------------------------------------------------------------
\makeatletter
\newcommand{\coursecode}[1]{\gdef\@coursecodetext{#1}}
\gdef\@coursecodetext{}
\makeatother

% -----------------------------------------------------------------------------
% Beamer theme assignments (only apply under Beamer).
%
% We detect Beamer via \@ifclassloaded{beamer}, which is the documented,
% non-internal way. This must be wrapped in \makeatletter/\makeatother
% because \@ifclassloaded contains an @-catcode character.
% -----------------------------------------------------------------------------
\makeatletter
\@ifclassloaded{beamer}{%
  \usefonttheme{professionalfonts}
  \setbeamercolor{structure}{fg=primary-blue}
  \setbeamercolor{frametitle}{fg=primary-blue}
  \setbeamercolor{title}{fg=primary-blue}
  \setbeamercolor{subtitle}{fg=primary-gold}
  \setbeamercolor{author}{fg=primary-blue}
  \setbeamercolor{institute}{fg=neutral}
  \setbeamercolor{date}{fg=primary-gold}
  \setbeamercolor{itemize item}{fg=primary-blue}
  \setbeamercolor{itemize subitem}{fg=primary-gold}
  \setbeamercolor{alerted text}{fg=primary-gold}
  \setbeamercolor{block title}{fg=white, bg=primary-blue}
  \setbeamercolor{block body}{bg=light-bg}
  % Semantic block variants: block=definition/concept (blue), exampleblock=
  % worked example (green/positive), alertblock=key takeaway or caution (gold).
  % Keep to <=2 colored boxes per slide across ALL three variants (INV-7).
  \setbeamercolor{block title example}{fg=white, bg=positive}
  \setbeamercolor{block body example}{bg=light-bg}
  \setbeamercolor{block title alerted}{fg=white, bg=primary-gold}
  \setbeamercolor{block body alerted}{bg=light-bg}

  % Minimal footer: course code (left) + page X / N (right), no navigation clutter
  \setbeamertemplate{navigation symbols}{}
  \setbeamertemplate{footline}{%
    \color{neutral}\footnotesize\hspace{0.7em}\@coursecodetext\hfill
    \insertframenumber/\inserttotalframenumber\hspace{0.7em}\vspace{0.3em}}
}{}
\makeatother

% -----------------------------------------------------------------------------
% TikZ setup — libraries the snippet gallery uses
% -----------------------------------------------------------------------------
\usepackage{tikz}
\usetikzlibrary{arrows.meta, positioning, calc, decorations.pathreplacing,
                fit, shapes.geometric, backgrounds}

% Shared TikZ styles — snippets and hand-written diagrams can pick these up
% via `\begin{tikzpicture}[every node/.style={...}]` or by naming specific
% styles (e.g., `\node[dag-node] (X) at (0,0) {X};`).
\tikzset{
  % Node styles for causal DAGs and similar
  dag-node/.style = {
    draw=primary-blue, fill=light-bg, rounded corners=2pt,
    minimum width=1.2cm, minimum height=0.9cm, align=center, font=\small
  },
  decision-node/.style = {
    draw=primary-gold, fill=white, diamond, aspect=1.8,
    minimum width=1.8cm, minimum height=1.2cm, align=center, font=\small,
    inner sep=1pt
  },
  % Edge styles — observed vs counterfactual
  observed-edge/.style = {-{Latex[length=2.5mm]}, thick, draw=jet},
  counterfactual-edge/.style = {-{Latex[length=2.5mm]}, thick, draw=neutral, dashed},
  confound-edge/.style = {-{Latex[length=2.5mm]}, thick, draw=negative, dashed},
  % Data-point styles for plots
  observed-dot/.style = {fill=primary-blue, draw=primary-blue, circle, inner sep=1.2pt},
  counterfactual-dot/.style = {fill=white, draw=neutral, circle, inner sep=1.2pt}
}

% -----------------------------------------------------------------------------
% Convenience macros
% -----------------------------------------------------------------------------
\newcommand{\muted}[1]{\textcolor{neutral}{#1}}
\newcommand{\key}[1]{\textcolor{primary-gold}{\textbf{#1}}}
\newcommand{\good}[1]{\textcolor{positive}{#1}}
\newcommand{\bad}[1]{\textcolor{negative}{#1}}

% Standout frame macro: full-bleed dark background for section transitions.
% Beamer-only (uses \begin{frame}/\setbeamercolor/\usebeamerfont) -- do not
% call from an article-class file (e.g. Notes/<CODE>/*-notes.tex). \muted,
% \key, \good, \bad, and \coursecode above are plain xcolor/gdef and are
% safe to use from either class; verified when Notes/article-class support
% was added.
% Expand: \transitionslide{Your transition title}
\newcommand{\transitionslide}[1]{%
  \begingroup
  \setbeamercolor{background canvas}{bg=primary-blue}
  \begin{frame}[plain]
    \centering
    \vfill
    {\usebeamerfont{title}\color{white}\Large #1\par}
    \vfill
  \end{frame}
  \endgroup
}

% Section divider frame (Beamer-only), an alternative to \transitionslide with a
% darker two-line style: near-black (jet) background, a small white label line
% above a larger gold title, and a thin gold rule along the bottom edge.
% Expand: \sectiondivider{Section 2}{Pointers}
\newcommand{\sectiondivider}[2]{%
  \begingroup
  \setbeamercolor{background canvas}{bg=jet}
  \begin{frame}[plain]
    \vspace*{3.0cm}
    {\color{white}\ttfamily\large #1\par}
    \vspace{0.5cm}
    {\color{primary-gold}\LARGE #2\par}
    \begin{tikzpicture}[remember picture, overlay]
      \draw[primary-gold, line width=2.5pt]
        ($(current page.south west)+(0,0.1)$) -- ($(current page.south east)+(0,0.1)$);
    \end{tikzpicture}
  \end{frame}
  \endgroup
}

% -----------------------------------------------------------------------------
% Channel watermark — "Concepts That Click" (YouTube).
%
% Article-class files (CompetitiveExam Q/A sets, Notes, Assignments) get a
% light diagonal draft watermark on every page plus a centered footer naming
% the channel. Beamer decks get a subtle TikZ background watermark instead
% (the footline already carries the course code). Centralized here so every
% MCQ PDF is branded without per-file edits.
% -----------------------------------------------------------------------------
\makeatletter
\@ifclassloaded{article}{%
  \usepackage{draftwatermark}
  \SetWatermarkText{Concepts That Click}
  \SetWatermarkScale{1}
  \SetWatermarkFontSize{2cm}
  \SetWatermarkLightness{0.86}
  \SetWatermarkAngle{45}
  \usepackage{fancyhdr}
  \pagestyle{fancy}
  \fancyhf{}
  \fancyfoot[C]{\footnotesize\textcolor{neutral}{Concepts That Click~|~YouTube: \href{https://www.youtube.com/@concepts.thatclick}{@concepts.thatclick}}}
  \renewcommand{\headrulewidth}{0pt}
  \renewcommand{\footrulewidth}{0pt}
  \fancypagestyle{plain}{%
    \fancyhf{}%
    \fancyfoot[C]{\footnotesize\textcolor{neutral}{Concepts That Click~|~YouTube: \href{https://www.youtube.com/@concepts.thatclick}{@concepts.thatclick}}}%
    \renewcommand{\headrulewidth}{0pt}%
    \renewcommand{\footrulewidth}{0pt}%
  }
}{}
\@ifclassloaded{beamer}{%
  \setbeamertemplate{background}{%
    \begin{tikzpicture}[remember picture, overlay]
      \node[rotate=45, opacity=0.055, align=center]
        at (current page.center)
        {\Huge\textcolor{neutral}{Concepts That Click}};
    \end{tikzpicture}%
  }
}{}
\makeatother 
\coursecode{JKSSB}
\title{Lesson 09 Practice: Virtual Memory}
\author{JKSSB Computer Awareness}
\date{\today}
\begin{document}
\maketitle
\noindent\textbf{Provenance:} All 20 items are original JKSSB-pattern
questions written for this lesson. None is a real PYQ; each is labelled
\textit{[Original]}.

\begin{enumerate}[label=\textbf{Q\arabic*.}, leftmargin=*]
\item \textit{[Original]} Virtual memory is best described as:
  \begin{enumerate}[label=\Alph*.]
    \item extra physical RAM added to the computer
    \item a technique that uses secondary storage to extend the addressable
      memory beyond physical RAM
    \item a type of cache built into the CPU
    \item a read-only memory used to store the BIOS
  \end{enumerate}
\item \textit{[Original]} Compared with RAM, virtual memory is:
  \begin{enumerate}[label=\Alph*.]
    \item faster
    \item exactly the same speed
    \item slower
    \item always disabled
  \end{enumerate}
\item \textit{[Original]} Which statement about virtual memory is correct?
  \begin{enumerate}[label=\Alph*.]
    \item It is additional physical RAM.
    \item It is disk-backed and slower than RAM.
    \item It is faster than cache.
    \item It removes the need for RAM.
  \end{enumerate}
\item \textit{[Original]} The main reason virtual memory exists is to:
  \begin{enumerate}[label=\Alph*.]
    \item increase the clock speed of the processor
    \item allow programs and data larger than physical RAM to run, while
      giving each process a large contiguous view of memory
    \item replace the CPU with a faster unit
    \item store the firmware permanently
  \end{enumerate}
\item \textit{[Original]} Expand MMU.
  \begin{enumerate}[label=\Alph*.]
    \item Memory Management Unit
    \item Main Memory Utility
    \item Multiprocessor Memory Unit
    \item Machine Mapping Unit
  \end{enumerate}
\item \textit{[Original]} The MMU is responsible for:
  \begin{enumerate}[label=\Alph*.]
    \item translating virtual addresses into physical addresses
    \item storing files permanently on disk
    \item cooling the processor
    \item displaying output on the monitor
  \end{enumerate}
\item \textit{[Original]} In paging, the fixed-size blocks into which a
  program's virtual address space is divided are called:
  \begin{enumerate}[label=\Alph*.]
    \item segments
    \item pages
    \item frames
    \item sectors
  \end{enumerate}
\item \textit{[Original]} In paging, the fixed-size blocks of physical memory
  into which pages are loaded are called:
  \begin{enumerate}[label=\Alph*.]
    \item pages
    \item frames
    \item segments
    \item registers
  \end{enumerate}
\item \textit{[Original]} Which statement correctly distinguishes paging from
  segmentation?
  \begin{enumerate}[label=\Alph*.]
    \item Paging uses variable-size segments; segmentation uses fixed-size
      pages.
    \item Paging divides memory into fixed-size pages and frames; segmentation
      divides a program into variable-size logical segments.
    \item Both paging and segmentation divide memory into equal fixed-size
      blocks.
    \item Segmentation divides memory into equal frames; paging uses
      variable-size logical units.
  \end{enumerate}
\item \textit{[Original]} Segmentation divides a program into:
  \begin{enumerate}[label=\Alph*.]
    \item equal fixed-size pages
    \item variable-size logical segments
    \item physical frames of equal size
    \item cache lines
  \end{enumerate}
\item \textit{[Original]} A page fault occurs when:
  \begin{enumerate}[label=\Alph*.]
    \item a hardware component develops a physical fault
    \item the referenced page is not currently in RAM and must be loaded from
      disk
    \item the contents of RAM become corrupted
    \item the processor overheats
  \end{enumerate}
\item \textit{[Original]} Demand paging loads a page into memory:
  \begin{enumerate}[label=\Alph*.]
    \item only when it is referenced and not already present
    \item only if the whole program is loaded at startup
    \item never, because pages stay on disk
    \item only during the boot process
  \end{enumerate}
\item \textit{[Original]} Which statement about virtual memory is \textbf{NOT}
  correct?
  \begin{enumerate}[label=\Alph*.]
    \item Virtual memory uses secondary storage.
    \item Virtual memory is slower than RAM.
    \item Virtual memory is additional physical RAM.
    \item Virtual memory lets a process use a large contiguous address space.
  \end{enumerate}
\item \textit{[Original]} Evaluate the statements:
  \begin{enumerate}[label=\Roman*.]
    \item Virtual memory uses secondary storage to extend memory.
    \item Virtual memory is slower than RAM.
    \item Virtual memory is extra physical RAM.
  \end{enumerate}
  \begin{enumerate}[label=\Alph*.]
    \item I and II only
    \item I only
    \item II and III only
    \item I, II, and III
  \end{enumerate}
\item \textit{[Original]} Swapping refers to:
  \begin{enumerate}[label=\Alph*.]
    \item moving data or pages between RAM and disk (secondary storage)
    \item exchanging the CPU for a GPU
    \item changing the monitor resolution
    \item moving data from cache to a CPU register
  \end{enumerate}
\item \textit{[Original]} Thrashing is best described as:
  \begin{enumerate}[label=\Alph*.]
    \item excessive paging that leaves little time for useful work
    \item a fast cache hit
    \item a type of segmentation
    \item the automatic addition of extra RAM
  \end{enumerate}
\item \textit{[Original]} Which statement correctly distinguishes cache from
  virtual memory?
  \begin{enumerate}[label=\Alph*.]
    \item Cache is disk-backed; virtual memory sits on the CPU chip.
    \item Cache is small, fast memory close to the CPU; virtual memory uses
      slower, disk-backed storage.
    \item Both cache and virtual memory are forms of secondary storage.
    \item Cache extends RAM using disk; virtual memory speeds up the CPU.
  \end{enumerate}
\item \textit{[Original]} The office desktop (EX-01) has 16\,GB of RAM and
  runs an application whose working set is larger than RAM. What does the
  operating system do?
  \begin{enumerate}[label=\Alph*.]
    \item It refuses to run the application.
    \item It uses virtual memory, keeping some pages on disk and loading them
      on demand.
    \item It converts part of the disk into extra RAM chips.
    \item It raises the clock speed to compensate.
  \end{enumerate}
\item \textit{[Original]} From the perspective of the CPU, which of these is
  the slowest to access?
  \begin{enumerate}[label=\Alph*.]
    \item L1 cache
    \item RAM
    \item disk-backed virtual memory
    \item a CPU register
  \end{enumerate}
\item \textit{[Original]} Which pair is correctly matched?
  \begin{enumerate}[label=\Alph*.]
    \item Paging --- fixed-size pages and frames
    \item Segmentation --- fixed-size pages and frames
    \item Paging --- variable-size logical segments
    \item MMU --- permanent file storage on disk
  \end{enumerate}
\end{enumerate}
\end{document}
